\section{The declared retained-label packet}
\label{sec:tpc53-literal-packet}

\subsection{Rows, scales, and the static mask}

Fix \(h_0\ne0\), put \(H=\rad|h_0|\), and take one
source-compatible opened-\(d\) block.  The source and opposite rows are
\[
 m=\ell_m d_m,\qquad m_\gamma=\ell_\gamma d_\gamma,
 \qquad \ell_m,\ell_\gamma\in[L,2L],\quad
 d_m,d_\gamma\in[D,2D],
\]
where the \(\ell\)'s are prime, the \(d\)'s are squarefree and coprime
to \(H\).  Endpoint atoms may be assigned to either adjacent dyadic
cell; every compact box selected below lies strictly inside these
intervals.  The inherited size separation makes each integer row
determine its pair \((\ell,d)\) uniquely
\citep{WangTPC25,WangTPC36}.  Write
\begin{equation}
 Q=LD,\qquad QJ\asymp X,\qquad J=K/D,
 \qquad
 c_X=\frac{LDJ}{X},\quad c'_X=\frac{DJ}{K}=1.
 \label{eq:tpc53-scales}
\end{equation}
The literal static mask is
\begin{equation}
 K_{m\gamma}
 =\one_{\ell_m\ne\ell_\gamma}
  \one_{|m-m_\gamma|>\Delta}
  \one_{(d_m,d_\gamma)\le G},
 \qquad
 \Delta=QX^{-\kappa_0},\quad G=X^{\kappa_0}.
 \label{eq:tpc53-literal-mask}
\end{equation}
It remains literal throughout this paper.

Let \(H_0=O_{h_0}(1)\) denote the fixed orbit modulus and let
\(q=O_{h_0}(1)\) be a common multiple of every fixed row and orbit
residue modulus.  Fix one locally admissible source/opposite row-residue
group modulo \(q\) and one orbit class \(a\pmod {H_0}\).  Direct
splitting makes the periodic multiplier a constant \(\xi_a\) on this
group and class.  Cells with \(\xi_a=0\) are discarded.  No Fourier
normalization is attached to a direct cell.

The opposite-row universe is complete within its fixed dyadic and
residue supports.  Writing \(V\) for the inherited hard divisor cutoff,
we assume the nondegenerate interior conditions
\begin{equation}
 L>4D,\qquad 2D\le V,\qquad L,D\longrightarrow\infty,
 \label{eq:tpc53-nondegenerate-row-cell}
\end{equation}
and no independent opposite-row pruning.  Together with the fixed
positive-density unary and admissible residue support declared below,
these conditions put the packet in the TPC--50 nondegenerate interior
regime, where the coefficient-free opposite factor is a complete
prime--squarefree cell and the literal mask survives \citep{WangTPC50}.

\subsection{The representative four-factor profile}

The source coefficient is retained without alteration:
\begin{equation}
 a_m:=\gamma_m^{(1)}
 =\mu(d_m)(\log\ell_m)\omega_D(d_m)
   \psi_L(\ell_m/L)\zeta_m^{(1)},
 \qquad |\zeta_m^{(1)}|\ll1.
 \label{eq:tpc53-source-coefficient}
\end{equation}
In particular, \(a_m\) is independent of the pre-terminal cofactor
label \(r\).  The coefficient-free opposite row still has its fixed
unary shape.  We write it exactly as
\begin{equation}
 v_0(\gamma)=\eta_\gamma
 R\!\left(\log\frac{\ell_\gamma}{L},
          \log\frac{d_\gamma}{D}\right),\qquad
 |\eta_\gamma|=1,\qquad
 R\in C^\infty([0,\log2]^2),
 \label{eq:tpc53-opposite-unary-shape}
\end{equation}
on its nonzero direct support.  Here \(R\) is the product of the fixed
dyadic shapes and \(\eta_\gamma\) contains the M\"obius sign and any
remaining unit-modulus direct-cell phase.  The exponent \(e=0\) removes
the opposite prime logarithm; it does not replace \(R\) by one.

We now name the smooth product that was abbreviated by
\(\mathcal W\) in later interfaces.  For fixed
\(W_1,W_2,\Psi_1,\Psi_2\in C_c^\infty((0,\infty))\), put
\begin{align}
 \mathcal W_{m,\gamma}(j/J)
 &:={}
 W_1(mj/X)\Psi_1(d_mj/K)
 W_2(m_\gamma j/X)\Psi_2(d_\gamma j/K).
 \label{eq:tpc53-four-factor-profile}
\end{align}
This is the representative product obtained by placing one source-row
factor and one opposite-row factor into the actual opened-row
provenance \citep{WangTPC25}.  Each has the form
\(W(mj/X)\Psi(dj/K)\).  The same argument works verbatim for any fixed
finite product; four factors keep the allocation visible.

Introduce normalized variables
\begin{equation}
 \beta_m=(x_m,y_m,c_X,c'_X),\quad
 x_m=\log(\ell_m/L),\quad y_m=\log(d_m/D),
 \qquad
 q_\gamma=(x_\gamma,y_\gamma),\quad t_j=j/J.
 \label{eq:tpc53-normalized-variables}
\end{equation}
Then the four arguments in \eqref{eq:tpc53-four-factor-profile} are
\begin{equation}
 \begin{aligned}
 z_1(\beta,t)&=c_X\e^{x_m+y_m}t,
 &\qquad z_2(\beta,t)&=c'_X\e^{y_m}t,\\
 z_3(\beta,q,t)&=c_X\e^{x_\gamma+y_\gamma}t,
 &\qquad z_4(\beta,q,t)&=c'_X\e^{y_\gamma}t.
 \end{aligned}
 \label{eq:tpc53-four-arguments}
\end{equation}
The normalized scale parameters range over a fixed compact set after a
fixed dyadic subdivision.

\subsection{The pre-terminal lift}

Let \(\cI_X\) contain the retained pairs \((m,r)\), let \(\cO_X\) be the
complete opposite-row cell, and let
\[
 \cJ_{X,a}=\{j:j/J\in I_0, j\equiv a\pmod {H_0}\},
\]
where \(I_0\Subset(0,\infty)\) is one fixed interval containing the
possible orbit support.  Define
\begin{equation}
 \cH_X^{\rm lift}
 :=\bigoplus_{(m,r)\in\cI_X}
   \ell^2(\cJ_{X,a};\ell^2(\cO_X)).
 \label{eq:tpc53-lifted-space}
\end{equation}
For each outer summand, the map
\(e_\gamma\mapsto e_{\gamma,d_mr}\) is an isometry onto the physical
coefficient subspace.  Thus suppressing the deterministic residual
coordinate here does not change the atomic normalization.
The declared coefficient vector is
\begin{equation}
 G_X^{L}(m,r,j)
 =\sum_{\gamma\in\cO_X}
 a_mK_{m\gamma}\eta_\gamma R(q_\gamma)\xi_a
 \mathcal W_{m,\gamma}(j/J)e_{\gamma}.
 \label{eq:tpc53-literal-profile-vector}
\end{equation}
Consequently
\begin{equation}
 \boxed{
 \|G_X^L\|_{\cH_X^{\rm lift}}^2
 =\sum_{m,r,j,\gamma}
 |a_m|^2K_{m\gamma}|R(q_\gamma)|^2|\xi_a|^2
 |\mathcal W_{m,\gamma}(j/J)|^2.}
 \label{eq:tpc53-exact-lifted-energy}
\end{equation}
Every sum in \eqref{eq:tpc53-exact-lifted-energy} is over the retained
sets just defined.

The label \(r\) is repeated orthogonally only in
\eqref{eq:tpc53-lifted-space}.  The terminal residual vector later
multiplies by \(\mu(r)\Lambda(pr)\), sums coherently over \(p\), and
groups \(s=d_mr\) \citep{WangTPC48}.  Distinct \((m,r)\) can then
collide.  Thus \eqref{eq:tpc53-exact-lifted-energy} is not an equality
for that downstream norm.

This boundary can be written as an exact operator.  Put
\(\cR_X^{\rm res}:=\{d_mr:(m,r)\in\cI_X\}\), let
\(\Phi_0\) be the fixed inherited smooth terminal cutoff, and put
\(\Lambda(n)=(\log n)\Phi_0(n/X)\),
and define
\[
 \mathcal T_{h_0}:\cH_X^{\rm lift}\longrightarrow
 \ell^2(\cO_X\times\cR_X^{\rm res}\times\cJ_{X,a})
\]
on an arbitrary lifted array \(g=(g_{m,r,j,\gamma})\) by
\begin{equation}
 (\mathcal T_{h_0}g)_{\gamma,s,j}
 :=
 \sum_{\substack{m,r,p\\s=d_mr\\pr-mj=h_0}}
 \mu(r)\Lambda(pr)g_{m,r,j,\gamma}.
 \label{eq:tpc53-terminal-grouping-map}
\end{equation}
Here \(p\) is prime and every variable in the sum is restricted to the
inherited finite terminal, row, cofactor, and orbit supports.  With
those supports understood, the equation \(pr-mj=h_0\) determines the
admissible terminal atoms.
The TPC--48 physical left residual output is
\(\mathcal T_{h_0}G_X^L\).  Its later identification with the original
large-prime residual carries one inherited global sign.  Its norm is
\begin{equation}
 \|\mathcal T_{h_0}G_X^L\|^2
 =
 \sum_{\gamma,s,j}
 \left|
 \sum_{\substack{m,r,p\\s=d_mr\\pr-mj=h_0}}
 \mu(r)\Lambda(pr)[G_X^L(m,r,j)]_\gamma
 \right|^2,
 \label{eq:tpc53-terminal-cross-term-norm}
\end{equation}
and therefore contains cross terms absent from
\eqref{eq:tpc53-exact-lifted-energy}.

Cauchy--Schwarz on each output fiber gives the upper bound
\begin{equation}
 \|\mathcal T_{h_0}g\|^2
 \le
 \left(
  \sup_{\gamma,s,j}
  \sum_{\substack{m,r,p\\s=d_mr\\pr-mj=h_0}}
  |\mu(r)\Lambda(pr)|^2
 \right)
 \|g\|_{\cH_X^{\rm lift}}^2.
 \label{eq:tpc53-terminal-upper-bound}
\end{equation}
No corresponding lower bound follows from multiplicity alone: every
output fiber containing two input atoms has a nontrivial kernel as a
linear functional on unrestricted arrays.  This is a logical
non-implication, not a claim that the distinguished arithmetic array
\(G_X^L\) belongs to that kernel.  Moreover, the residual output itself
belongs to the low-sign-averaged Gram route; its cross-prime inner
products have not been identified with the native affine physical Gram
\citep{WangTPC48}.
